\documentclass[10pt,a4paper]{amsart}
\usepackage{amsmath,amssymb,amsthm,graphicx}
% Retained preamble for the Stacks Project chapter snapshot in this directory.
% The chapter's own preamble.tex is a copy of the shared Stacks Project preamble; the block
% below is sliced MECHANICALLY, line for line, from it. Nothing was added, removed or
% redefined; the wrapper class, the xr-hyper/hyperref wiring and the externaldocument list
% are omitted because the assembler supplies the document class.
%
% lines taken from preamble.tex: 3, 5, 9, 11, 12, 156, 157, 158, ... (51 lines total)
%
% package coverage restored from the chapter's own preamble (2026-10-13): xy, xy-2cell

\usepackage[all]{xy}
\xyoption{2cell}
\UseAllTwocells

\newenvironment{reference}{\comment}{\endcomment}
\newenvironment{slogan}{\comment}{\endcomment}
\newenvironment{history}{\comment}{\endcomment}
\theoremstyle{plain}
\newtheorem{theorem}[subsection]{Theorem}
\newtheorem{proposition}[subsection]{Proposition}
\newtheorem{lemma}[subsection]{Lemma}
\theoremstyle{definition}
\newtheorem{definition}[subsection]{Definition}
\newtheorem{example}[subsection]{Example}
\newtheorem{exercise}[subsection]{Exercise}
\newtheorem{situation}[subsection]{Situation}
\theoremstyle{remark}
\newtheorem{remark}[subsection]{Remark}
\newtheorem{remarks}[subsection]{Remarks}
\numberwithin{equation}{subsection}
\def\lim{\mathop{\mathrm{lim}}\nolimits}
\def\colim{\mathop{\mathrm{colim}}\nolimits}
\def\Spec{\mathop{\mathrm{Spec}}}
\def\Hom{\mathop{\mathrm{Hom}}\nolimits}
\def\Ext{\mathop{\mathrm{Ext}}\nolimits}
\def\SheafHom{\mathop{\mathcal{H}\!\mathit{om}}\nolimits}
\def\SheafExt{\mathop{\mathcal{E}\!\mathit{xt}}\nolimits}
\def\Sch{\mathit{Sch}}
\def\Mor{\mathop{\mathrm{Mor}}\nolimits}
\def\Ob{\mathop{\mathrm{Ob}}\nolimits}
\def\Sh{\mathop{\mathit{Sh}}\nolimits}
\def\NL{\mathop{N\!L}\nolimits}
\def\CH{\mathop{\mathrm{CH}}\nolimits}
\def\proetale{{pro\text{-}\acute{e}tale}}
\def\etale{{\acute{e}tale}}
\def\QCoh{\mathit{QCoh}}
\def\Ker{\mathop{\mathrm{Ker}}}
\def\Im{\mathop{\mathrm{Im}}}
\def\Coker{\mathop{\mathrm{Coker}}}
\def\Coim{\mathop{\mathrm{Coim}}}
\DeclareMathSymbol{\boxtimes}{\mathbin}{AMSa}{"02}
\def\QCohstack{\mathcal{QC}\!\mathit{oh}}
\def\Cohstack{\mathcal{C}\!\mathit{oh}}
\def\Spacesstack{\mathcal{S}\!\mathit{paces}}
\def\Quotfunctor{\mathrm{Quot}}
\def\Hilbfunctor{\mathrm{Hilb}}
\def\Curvesstack{\mathcal{C}\!\mathit{urves}}
\def\Polarizedstack{\mathcal{P}\!\mathit{olarized}}
\def\Complexesstack{\mathcal{C}\!\mathit{omplexes}}
\def\Pic{\mathop{\mathrm{Pic}}\nolimits}
\def\Picardstack{\mathcal{P}\!\mathit{ic}}
\def\Picardfunctor{\mathrm{Pic}}
\def\Deformationcategory{\mathcal{D}\!\mathit{ef}}
\usepackage{cleveref}
\crefname{originallemma}{original lemma}{original lemmas}
\Crefname{originallemma}{Original lemma}{Original lemmas}
\makeatletter
\newlabel{definition-hypercovering-variant}{{09VU}{source}{}{sourcepointer-1}{}}
\newlabel{definition-hypercovering-variant@cref}{{[originallemma][0][]09VU}{[1][0][]source}{}{}{}}
\makeatother
\crefname{originallemma}{original lemma}{original lemmas}
\Crefname{originallemma}{Original lemma}{Original lemmas}
\makeatletter
\newlabel{simplicial-remark-existence-cosk}{{09VS}{source}{}{sourcepointer-2}{}}
\newlabel{simplicial-remark-existence-cosk@cref}{{[originallemma][0][]09VS}{[1][0][]source}{}{}{}}
\makeatother
\crefname{originallemma}{original lemma}{original lemmas}
\Crefname{originallemma}{Original lemma}{Original lemmas}
\makeatletter
\newlabel{sites-lemma-pullback-representable-sheaf}{{04D3}{source}{}{sourcepointer-3}{}}
\newlabel{sites-lemma-pullback-representable-sheaf@cref}{{[originallemma][0][]04D3}{[1][0][]source}{}{}{}}
\makeatother
\title{Morphisms of sites preserve hypercoverings}
\author{The Stacks Project Authors}
\date{stacks-project:hypercovering.tex:local-snapshot-20261012; GNU FDL 1.2 or later}
\begin{document}
\maketitle
\noindent\textit{Editorial note.} Statement span is lines 3037-3045 of sources/stacks-hypercovering/pinned-original-source.tex; proof span is lines 3047-3078, the complete author proof (30 body lines).\par\medskip
\noindent\textit{Editorial note.} Navigation stubs: 3 cited-but-not-included labels are supplied as non-load-bearing pointers with their pinned tags where the snapshot tag index has them.\par\medskip
\noindent\textit{Editorial note.} extra\_packages is empty on purpose: the retained preamble already loads xy with options and a second declaration clashes in the assembler wrapper.\par\medskip
\noindent\textit{Editorial note.} Text is retained verbatim, including the author's own typos and placeholders.\par\medskip
\noindent\textit{Editorial note.} Queue index 15 of queue\_hypercovering.json (pb=1331, stmt 3037-3045, proof 3047-3078).\par\medskip
\noindent\textit{Editorial note.} stacks-hypercovering wave-1 (hc01): fresh adjudication accept (2026-10-14); pool per standing fallback order (hypercovering -> pic -> spaces-more-cohomology) after md01 closed the moduli pool; hypercovering.tex 53 labels/34 with proof judged 26 accept / 8 reject / 14 pb=0 env-dup rows excluded; spans cite-free and reference/slogan/history-free so no patchenv needed.\par\medskip
\noindent\textit{Allowed standard prerequisites (editorial).}\begin{itemize}
\item Explicit prerequisite boundary. The cited inputs are prerequisites; the argument in the retained span is not replaced by any cited result.
\end{itemize}\medskip
\ifdefined\hypertarget\hypertarget{sourcepointer-1}{}\fi
\noindent\textit{Source reference.} Original lemma 09VU: sources/stacks-hypercovering/pinned-original-source.tex, source lines 1285--1298. Cited by the retained statement or proof but not part of the collected spans; supplied only as a navigation pointer. The argument does not depend on its content.\par\medskip
\ifdefined\hypertarget\hypertarget{sourcepointer-2}{}\fi
\noindent\textit{Source reference.} Original lemma 09VS: sources/stacks-hypercovering/pinned-original-source.tex, source lines 3078--3078. Cited by the retained statement or proof but not part of the collected spans; supplied only as a navigation pointer. The argument does not depend on its content.\par\medskip
\ifdefined\hypertarget\hypertarget{sourcepointer-3}{}\fi
\noindent\textit{Source reference.} Original lemma 04D3: sources/stacks-hypercovering/pinned-original-source.tex, source lines 3078--3078. Cited by the retained statement or proof but not part of the collected spans; supplied only as a navigation pointer. The argument does not depend on its content.\par\medskip
\begin{lemma}
\label{lemma-hypercovering-morphism-sites}
Let $f : \mathcal{C} \to \mathcal{D}$ be a morphism of sites
given by the functor $u : \mathcal{D} \to \mathcal{C}$.
Assume $\mathcal{D}$ and $\mathcal{C}$ have equalizers and
fibre products and $u$ commutes with them.
If a simplicial object $K$ of $\text{SR}(\mathcal{D})$
is a hypercovering, then $u(K)$ is a hypercovering.
\end{lemma}

\begin{proof}
If we write $K_n = \{U_{n, i}\}_{i \in I_n}$ as in the introduction
to this section, then $u(K)$ is the object of $\text{SR}(\mathcal{C})$
given by $u(K_n) = \{u(U_i)\}_{i \in I_n}$.
By Sites, Lemma \ref{sites-lemma-pullback-representable-sheaf}
we have $f^{-1}h_U^\# = h_{u(U)}^\#$ for $U \in \Ob(\mathcal{D})$.
This means that $f^{-1}F(K_n)^\# = F(u(K_n))^\#$ for all $n$.
Let us check the conditions (1), (2), (3) for $u(K)$ to be a
hypercovering from Definition \ref{definition-hypercovering-variant}.
Since $f^{-1}$ is an exact functor, we find that
$$
F(u(K_0))^\# = f^{-1}F(K_0)^\# \to f^{-1}* = *
$$
is surjective as a pullback of a surjective map and we get (1).
Similarly,
$$
F(u(K_1))^\# = f^{-1}F(K_1)^\# \to
f^{-1} (F(K_0) \times F(K_0))^\# = F(u(K_0))^\# \times F(u(K_0))^\#
$$
is surjective as a pullback and we get (2). For condition (3),
in order to conclude by the same method it suffices if
$$
F((\text{cosk}_n \text{sk}_n u(K))_{n + 1})^\# =
f^{-1}F((\text{cosk}_n \text{sk}_n K)_{n + 1})^\#
$$
The above shows that $f^{-1}F(-) = F(u(-))$. Thus it suffices to show
that $u$ commutes with the limits used in defining
$(\text{cosk}_n \text{sk}_n K)_{n + 1}$ for $n \geq 1$.
By Simplicial, Remark \ref{simplicial-remark-existence-cosk}
these limits are finite connected limits and $u$ commutes with these
by assumption.
\end{proof}

\end{document}
